\documentclass[11pt]{article}
\usepackage[utf8]{inputenc}
\usepackage[margin=1.15in]{geometry}
\usepackage{microtype}
\usepackage{booktabs}
\usepackage{array}
\usepackage{amsmath,amssymb,amsthm}
\usepackage[hidelinks]{hyperref}

\newtheorem{proposition}{Proposition}
\theoremstyle{remark}
\newtheorem{remark}{Remark}

\title{On Mental Imagery}
\author{Flyxion\\Independent Researcher}
\date{October 2026}

\begin{document}
\maketitle

\begin{abstract}
The vocabulary of seeing, picturing, and imagining pervades the description of thought, and it has repeatedly been taken as evidence that thought is carried out in internally experienced pictures. This essay follows that inference through a sequence of cases: Galton's breakfast-table inquiry of 1880 and its later reanalysis, Swedenborg's waking visions, Temple Grandin's account of thinking in pictures, Aldous Huxley's self-description as a poor visualizer, and Kant's theory of schematism, in which imagination is a rule-governed capacity and not a gallery of pictures. Frege and the later Wittgenstein supply the philosophical analysis of why an accompanying image cannot constitute understanding, and evidence from congenital blindness and from aphantasia supplies a test of what visual language requires of those who use it. The essay separates three things that the word \emph{image} tends to run together: a representation that is available to a system, the conscious experience of something like seeing, and the competence to perform tasks associated with visualization. It does not argue that mental imagery is illusory. It argues that a practice shared by an entire linguistic community is weak evidence about the uniformity of the experience its words appear to describe.
\end{abstract}

\section{Three Things the Word Image Covers}

Ordinary speech about the mind is saturated with the vocabulary of vision. A person pictures a future, sees a point, takes a view, brings a difficulty into focus, and forms an idea, a word that descends from the Greek verb for seeing.\footnote{The Greek \emph{idea} and \emph{eidos} are related to \emph{idein}, to see; the Latin \emph{videre} gives \emph{evident}, \emph{provide}, and \emph{vision}. Etymology is cited here as history of the vocabulary, not as an argument about present meaning, a distinction developed in Section 9.} The vocabulary is so pervasive that it invites a theory, according to which thinking consists, at least in its paradigm cases, in the inspection of internal pictures, and the language of vision is simply the literal description of what thinkers undergo. The theory has a long philosophical pedigree and an equally long record of embarrassment, because whenever the question has been put directly to people, they have disagreed about whether anything of the kind occurs in them.

The phrase \emph{mental image} bundles together at least three distinct claims, and much of the confusion surrounding it arises from sliding between them. The first concerns \emph{representation}: that some state of a person or system carries information about a scene, an object, or a configuration in a form that can be used in later tasks. The second concerns \emph{experience}: that the person undergoes something quasi-perceptual, a presentation to consciousness with some of the character of seeing. The third concerns \emph{competence}: that the person can describe, draw, recognize, rotate, compare, or otherwise reason about the scene in the ways that visualizing is expected to support. These will be written $R$, $E$, and $C$ for brevity, and the argument of the essay is that they are separable. A person can display $C$ with uncertain or absent $E$; a person can report vivid $E$ whose relation to $R$ and to $C$ is a further empirical matter; and the words used to describe all three are drawn from a common stock whose grammar does not settle which is meant.\footnote{The labels are expository. No claim is made that $R$, $E$, and $C$ exhaust the senses of \emph{image}, only that these three are routinely run together and that each can be questioned separately.}

The essay proceeds historically and then analytically. Sections 2 to 6 consider five figures who stand in different relations to the question: Galton, who found that people differ; Swedenborg and Grandin, who report experiences of great visual richness in very different settings; Huxley, who reports the opposite; and Kant, who offers a theory of imagination in which pictures are not the central item. Section 7 recalls the earlier dispute over imageless thought. Sections 8 and 9 turn to the philosophy of language, where the question is why the vocabulary of vision should be expected to track the phenomenology of vision at all. Section 10 considers the evidence from congenital blindness, Section 11 the methodological problem of inferring mechanisms from introspective reports, and Section 12 the relation between imagery and intellectual competence. The conclusion states what the argument does and does not establish.

A note on scope is required at the outset. The thesis is not that mental imagery is a fiction, nor that those who report vivid imagery are mistaken about their own experience. The reports of vivid imagery are evidence, and are treated as evidence throughout. The thesis concerns the step from the existence of a widely shared way of speaking to the existence of a widely shared kind of experience, and it holds that this step is not licensed.

\section{Galton and the Discovery of Divergent Inner Experience}

Francis Galton's \emph{Statistics of Mental Imagery}, published in \emph{Mind} in 1880, is the usual starting point for the empirical study of individual differences in visualization, and it deserves a careful reading because both its design and its later reception have shaped the field. The instrument was a short schedule of questions. Respondents were asked to call to mind a familiar scene, for which the breakfast table of that morning served as the example, and then to describe the resulting picture under three headings: its illumination, the definition of its parts, and the naturalness of its colouring. Under illumination the question was whether the image was dim or fairly clear and whether its brightness matched the original; under definition, whether all the objects were well defined at once or only a narrower region was sharp, as compared with a real scene; under colouring, whether the colours were distinct and natural.

The memoir analyses the male returns. The adult sample comprised one hundred men, of whom at least half were distinguished in science or other intellectual work, nineteen of them Fellows of the Royal Society. A second sample consisted of schoolboys from a single school, the pupils of its science master, divided into an upper and a lower group so that each could serve as a check on the other. Replies had also been received from women, from children, and from members of the general public and of a French learned society, but these were not analysed in the memoir. Galton ranked the responses from highest to lowest and reported them by a small number of key positions, namely the middle, the quartiles, the octiles, and the extremes, which allowed him to describe the shape of the distribution without assigning numerical scores to experiences that have no natural unit.\footnote{On the sampling, the ranking by key positions, and the proportions reported for the boys, see Galton (1880); the memoir also reports on the apparent distance of images and on visualized numerals, which are not discussed here.}

The striking feature of the paper is the range of reports. At one end were respondents who described images of illusory perceptual clarity, in some cases so strong that they were unsure how to describe the difference from sight. At the other end were respondents who could recall the arrangement of the table, the number of people seated at it, and the items upon it, and who nevertheless denied having anything that could be called a picture. Galton relates that many of the scientific men to whom he first applied protested that mental imagery was unknown to them, and he compared their position to that of the colour-blind, who are unaware of what they lack. Some respondents, he reports, had taken the language of mental pictures to be a figure of speech, and one wrote that to describe recollection as seeing an image is only a way of talking. The point is hard to overstate. A population had been using a vocabulary of internal sight for centuries, and a questionnaire revealed that some of its members had always understood the vocabulary as figurative while others had understood it literally, with neither group having had any reason to suspect the difference.\footnote{Galton (1880) describes the colour-blind comparison and the figurative reading of \emph{mental picture}; both are paraphrased here and not quoted.}

The paper also contains a theory, and the theory has had a more troubled history than the data. Galton argued that an over-readiness to see clear mental pictures works against habits of highly generalized and abstract thought, that a faculty little used by hard thinkers is easily lost, and that the highest minds therefore probably keep visual imagery in a subordinate state, available when needed. He suggested that in scientists the weakness of visual imagery is made good by other modes of thought, and he linked these chiefly to the motor sense. This claim, that scientists as a class are deficient in visual imagery, was absorbed into the secondary literature and repeated for more than a century. A reanalysis by Brewer and Schommer-Aikins (2006), who replicated the breakfast-table inquiry with modern scientists and undergraduates and re-examined Galton's published data, found that Galton's own tables do not support the claim. In the 1880 data and in the modern replication alike, all groups reported substantial imagery on recollective tasks of this kind, no modern scientist was totally lacking in visual imagery, and very few reported feeble imagery. The authors conclude that Galton's generalization was an example of theory-laden interpretation, driven by the initial replies of a few very salient scientists who reported little or none.\footnote{Brewer and Schommer-Aikins (2006). The replication used a modern scientific sample and undergraduates and compared the results with Galton's published tables, which the authors argue do not support his generalization.}

The reanalysis is instructive for the present argument in two ways. It removes a tempting but unsupported moral, namely that abstract intellectual work is bought at the price of imagery, and it leaves untouched what the original paper established beyond reasonable doubt, which is that reports of imagery diverge, that the divergence extends in some individuals to the denial of anything picture-like, and that the divergence had gone unnoticed because the shared vocabulary concealed it. The surviving result is more modest than Galton's theory and more important. The discovery concerns the disagreement itself and not the distribution of imagery across professions.

Several limitations of the method should be stated. The sample was selected, composed largely of acquaintances and correspondents of the investigator, and it was entirely male in the analysed portion. The measure was introspective report, so that any two respondents may have applied different standards of clarity to the same inner state, or the same standard to different states. The questions presupposed that the respondent had a picture whose properties could be compared with those of a real scene, and a respondent without such a picture had to decide how far the question applied to him. The ranking procedure, which is sensible given the absence of a unit, also hides the question of whether the scale is one dimension or several. It should be added that Galton's wider programme was bound up with hereditarian and eugenic commitments that are rejected here; nothing in the argument of this essay depends on them, and the empirical observation about divergent reports can be separated from them without remainder.\footnote{The separation is the same one made throughout the essay between a report and the theory built on it: the divergence of responses is an observation, and the hereditarian use Galton made of imagery is an interpretation that the observation does not require.}

What Galton's inquiry shows, then, is a pattern of disagreement under shared vocabulary. Later work, from the vividness questionnaire of Marks (1973) onward, has refined the instruments, and the term \emph{aphantasia}, introduced by Zeman and colleagues in 2015 for the lifelong absence of voluntary visual imagery, has given the extreme end a name. Each of these developments rests on the same observation that Galton made accidentally, which is that people who use the same words about recollection may be reporting very different experiences, and may live for decades without learning of it.

\section{Swedenborg and the Clarity of Inner Vision}

Emanuel Swedenborg (1688 to 1772) offers a historical case at the opposite pole from Galton's non-visualizers, and he is a useful one because his testimony is abundant, consistent, and explicit about perceptual quality. Trained as a natural philosopher and engineer, and for years an assessor on the Swedish Board of Mines, he published extensively on mineralogy, anatomy, and cosmology before the middle of his life. A journal of dreams kept in 1743 and 1744 records vivid dream imagery and states between waking and sleep with unusual attention to their character. From 1745 onward he reported that his spiritual sight had been opened, and the large theological works of the following decades, among them the \emph{Arcana Caelestia} and \emph{Heaven and Hell}, are founded on accounts of what he saw and heard in the company of spirits and angels.\footnote{The dating of the opening of his spiritual sight to 1745, at the time of his residence in London, is conventional. The journal of dreams was not published in his lifetime. See the entry for Swedenborg in the references.}

Two features of these accounts bear on the question of imagery. The first is their insistence on perceptual clarity. Swedenborg repeatedly says that what he saw was seen as plainly as the objects of ordinary sight, and the works describe settings, buildings, gardens, garments, and the movements of persons in a degree of visual detail that is sustained across many volumes. The second is that these experiences are presented as occurring while he was awake, in full possession of his faculties, and as being at the same time distinct from ordinary perception. It would be implausible to treat this as an occasional use of picturesque language. The writings are, taken together, a long report of visually organized experience of great richness, and they are evidence about the phenomenology of the man who wrote them.

The proper use of this evidence needs to be stated with some precision. Contemporary hyperphantasia denotes unusually vivid imagery, typically that which a person can call up voluntarily, whereas Swedenborg's reports concern visions that he presents as received and not constructed, and that he interprets as encounters with realities independent of him. The distinctions among voluntary visualization, spontaneous imagery, dream, hypnagogic imagery, hallucination, and religious vision are real and cut across one another, and it would be a mistake to collapse the case into any one of them. What can be said is that Swedenborg's reports are strong testimony to exceptionally vivid visual experience, that its voluntariness is uncertain, and that the question of whether the objects he describes existed independently is a separate question from the one the essay asks.

That last separation matters for consistency. A researcher who accepts a modern respondent's statement that she can inspect a remembered kitchen with photographic clarity, as evidence that she has such experiences, should not reject Swedenborg's statements about what he saw solely on the ground that their content is religious. The testimony is admitted in both cases as evidence about the character of the experience, with the question of its causal origin held apart. Whether the experience is explained by an unusual capacity of imagery, by a pathological condition, by a trained practice of attention, or by the nature of the objects seen, is a question that these reports do not by themselves settle. Treating the reports in this way permits the historical record to be used without either endorsing the metaphysics or dismissing the experience.

The case also illustrates a historical problem, which is how a culture with an established vocabulary of visions, revelations, and apparitions distinguishes an unusually vivid imagination from the perception of another world. Kant addressed exactly this problem in 1766 in \emph{Dreams of a Spirit-Seer}, a work provoked by Swedenborg's claims. In the third chapter of its first part Kant offers a physiological conjecture, borrowed from the language of optics. In ordinary sensation, the mind locates an object at the point where the lines of the sense impression meet, and the images of imagination normally have this point inside the head, which is why a waking person recognizes them as fancies. If a disturbance of the nerves moves the point outside the head, the image is experienced as a real object in front of the person. Kant applies this to the confused and the deluded, and notes that sleepers and the drunk can undergo a similar displacement briefly. The conjecture was an attempt to explain how an image can present itself as a perception, and it takes for granted that such images exist in the person who reports them. It does not attempt to deny the experience, and in this respect it is closer to the position adopted here than a dismissal would have been.\footnote{Kant (1766), Part I, chapter 3. The tone of the work toward Swedenborg is ironic, and the physiological conjecture is offered as a hypothesis about how imagery could be mistaken for perception. The reference is to the passage as it appears in the 1899 English translation by Goerwitz and in the German original.}

\section{Grandin and Thinking in Pictures}

Swedenborg's imagery is placed within a metaphysical interpretation. A modern case of comparable reported vividness, placed within an entirely practical interpretation, is that of Temple Grandin, whose \emph{Thinking in Pictures} (1995, with an expanded edition in 2006) describes visual thought as the central medium of her cognition. She reports that she converts spoken and written words into pictures, that she can retrieve specific visual memories with great exactness, and that she can build and run equipment in her imagination, inspecting how a structure will behave under load or how an animal will move through a handling system before the design is built. In this account imagery is not a spectacle that accompanies thought; it is the working medium in which problems are solved, and the account is notable for tying the reported experience directly to engineering performance.\footnote{Grandin (1995; expanded edition 2006). The descriptions are self-report, and she herself separates her case from that of others on the autism spectrum.}

Grandin's description of concept formation is of particular interest to the philosophical history that follows. She reports that general concepts are built from collections of particular remembered instances: the concept of a kind of animal is, for her, a store of the individual animals she has seen, sorted by resemblance. This is very nearly the empiricist picture of abstract ideas that Locke, Berkeley, and Hume argued over, and it is a rare case in which a person claims to inhabit that picture rather than to theorize about it. The testimony is therefore valuable twice over, as a report of strong imagery and as a report of how such imagery might support general thought. The report also faces the problem that Kant and Berkeley raised, which is how a particular image can serve as an image of the general, and the account of sorting by resemblance leaves that problem in place, relocated to the sorting.

Grandin has also been candid about limits. She has written of her difficulty with algebra and of the contrast between her strengths in visual and spatial tasks and her weaknesses in abstract symbolic work, and she has described her discovery that other people's thought differs from hers, having at first assumed that everyone thinks in pictures. Her later writings go on to distinguish visual thinkers from pattern thinkers and verbal thinkers.\footnote{The three-way division appears in her later writing; it is a proposal of her own and has not been presented here as an established taxonomy.} The first of these observations is a direct counterexample to a simple equation of vivid imagery with general intellectual power, and the second repeats, in the register of personal experience, Galton's discovery that the assumption of uniformity is natural and false.

Work in cognitive psychology has given the contrast some structure. Kozhevnikov, Kosslyn, and Shephard (2005) distinguished object visualizers, whose imagery is pictorial, detailed, and vivid, from spatial visualizers, whose imagery is schematic and concerned with relations and transformations, and Blajenkova, Kozhevnikov, and Motes (2006) developed a questionnaire to measure the two dimensions separately.\footnote{Related work by the same group suggests that scientists and engineers tend toward the spatial dimension and visual artists toward the object dimension. This is a statement about group tendencies and is not used in the argument.} On this view the word \emph{imagery} covers more than one capacity, and the image-heavy account that Grandin gives corresponds most naturally to the object dimension, with transformation of objects in space as a further capacity. The distinction is consistent with the thesis of this essay: it shows that even among people who report abundant imagery, the experiences are not of one kind.

The Swedenborg and Grandin cases, taken together, show how the same type of report, a person saying that he or she sees things with perceptual clarity, can be embedded in radically different interpretive frameworks and attach to very different degrees of voluntary control. In one the imagery is received, in the other it is used. In neither case does the report by itself determine whether the imagery is representational in the sense of $R$, whether it is experienced in the sense of $E$, or whether it is the source of the competence $C$ that accompanies it. What the reports establish is that $E$ of great richness occurs.

\section{Huxley and the Impoverished Visual Imagination}

Where Swedenborg and Grandin supply testimony of abundance, Aldous Huxley supplies one of the best-known twentieth-century statements of scarcity. In \emph{The Doors of Perception} (1954) he describes himself as a poor visualizer and says that he has been one for as long as he can remember. Words, including the words of poets, do not ordinarily call up pictures in his mind. Memories do not return to him as scenes seen again, and the most that deliberate effort produces is a faint and fleeting image. The statement is autobiographical and explicit, it was made by a professional writer of great descriptive range, and it appears in a book devoted to altered perception, which gives it an unusual interest.

The report is not a report of complete aphantasia in the strict modern sense, because Huxley acknowledges some voluntary imagery, faint though it is. It is nevertheless a clear historical statement of markedly diminished visualization. Its further interest lies in what it shows about the drug experience he goes on to describe. A common expectation holds that psychedelic states consist of elaborate internally generated visual scenes, and the expectation is often taken as evidence for an inner theatre that the drug unlocks. Huxley's account suggests a different picture. The central transformations he reports concern present perception: colours intensified, objects acquiring an unusual intensity and significance, the folds of a garment or the legs of a chair becoming objects of absorbed attention. These are changes in how the visible world appears and not the production of imagery by a visually inventive mind. His report also does not claim a total absence of visual experience under the drug, and the contrast he draws is with his own baseline.\footnote{Huxley (1954). The mescaline session took place in 1953; the interpretation of the experience in terms of a reducing valve on a wider awareness is Huxley's own and is not endorsed here.} It is enough for present purposes that the ordinary and the altered states are described in terms of seeing, while the imagery capacity on which the theatre account depends was one he took himself to lack.

The comparison with Swedenborg is instructive and has a literary dimension. The title of Huxley's book is taken from William Blake's \emph{The Marriage of Heaven and Hell}, a work written in part against Swedenborg, whose \emph{Heaven and Hell} Blake had read and satirized. Huxley's own later volume on visionary experience carries a title that returns to the same pair.\footnote{Blake (1790) comments critically on Swedenborg in the \emph{Marriage}; Huxley (1956) treats visionary experience in the sequel to the earlier book.} The two writers share a topic and an interest in extraordinary experience, and they describe the relation of imagination to perception in nearly opposite terms. Swedenborg reports a surplus of visual presentation interpreted as perception of another world, and Huxley reports a poverty of voluntary imagery together with a transformed perception of this one. The difference should not be overinterpreted. It is a difference of two testimonies, not a measurement, and Huxley's well-documented visual impairment from youth is a reason for caution in drawing inferences about his imagery from his sensory history in either direction. What can be said is that a major writer on perception, expressing himself in some of the most vivid prose in the language, reports that the vividness of his writing did not rest on vivid inner pictures.

That final observation anticipates the later argument. A person can produce, comprehend, and value richly visual language without the experience that the language might be thought to describe. The expressive power of a text and the imagery of its author are two things, and the first is no guide to the second.

\section{Kant and the Imagination Without Pictures}

It is tempting to add Kant to this series as a presumptive non-visualizer, and the temptation should be named so that it can be resisted. No credible biographical evidence establishes that Kant lacked voluntary visual imagery, and the philosophical theory of imagination that he developed cannot be used as evidence about his own phenomenology, since a theory may be motivated by considerations that have nothing to do with the theorist's introspection. Reports from his students describe lively and picturesque lectures, which as the case of Huxley shows would be equally compatible with either condition.\footnote{The reports of his lectures come from reminiscences of students and are mentioned only to note that they decide nothing.} The conjecture that Kant was aphantasic therefore remains a conjecture, and the argument does not depend on it. What matters is the theory.

In the chapter on schematism in the \emph{Critique of Pure Reason}, Kant asks how the pure concepts of the understanding, which are not drawn from experience, can apply to the manifold of sensible intuition, which is. The mediating third thing must be homogeneous with both, and he identifies it with the schema, a product of the imagination in its productive rather than merely reproductive employment, determined in relation to time. The crucial point for present purposes is that a schema is explicitly distinguished from an image. An image is a particular, and a schema is a rule or procedure of the imagination for providing a concept with its image. The example Kant uses is the triangle. No single drawn or imagined triangle can be adequate to the concept of triangle in general, for the image would be right-angled or acute, equilateral or scalene, and the concept covers all of these.\footnote{Kant (1781/1787), the chapter on schematism, running from A137/B176 to A147/B187; the triangle passage falls in the middle of the chapter.} What is wanted is the procedure by which triangles of any form can be generated and recognized, and this is the schema. The same holds, he says, of empirical concepts: the concept of a dog is a rule by which the imagination can trace the general outline of a four-footed animal without being confined to any single shape drawn from experience.

The treatment has several consequences for the debate about imagery. First, it places imagination within cognition as an organizing operation and not as a faculty of inner display. Imagination synthesizes, and what it provides is the capacity to bring intuitions under concepts, which may or may not be accompanied by anything pictorial. Second, it identifies a problem that any picture theory of thought must face: a picture is particular, a thought is general, and an account is needed of how the one can serve as the other. Berkeley had pressed a version of this problem against Locke's abstract ideas, observing that he could not form an idea of a triangle that was neither oblique nor right-angled, and Locke's own wording of the triangle idea had been the occasion.\footnote{Berkeley (1710), Introduction; Locke (1689), Book IV. The passage in Locke is frequently cited as the target of the objection.} Descartes had earlier noticed that a figure of a thousand sides can be conceived distinctly and yet cannot be imagined distinctly, so that the two capacities come apart.\footnote{Descartes (1641), Sixth Meditation, the example of the chiliagon.} Kant's schematism is the most developed answer in that tradition, and it can be read as the claim that imagination in the sense that matters for knowledge is not the capacity to have pictures.

Third, the account is neutral on the empirical question of how people differ. A person with abundant visual imagery and a person with none can both possess the rule-governed capacity to apply the concept of a triangle, the former perhaps with visible illustration and the latter without. The theory gives no reason to expect uniformity of experience and good reason to expect uniformity of competence. For this reason the conjecture about Kant's own imagery is of secondary importance, and the stronger lesson is the philosophical one: a theory of imagination need not be a theory of internally visible images. It is also worth recording that the history of the schematism chapter has not been one of consensus. Kant described the operation as hidden in the depths of the human soul, and commentators have disagreed on what the schemata are. That disagreement does not weaken the point that Kant took the problem of the image and the general to be real.

\section{Imageless Thought and an Earlier Dispute}

The question whether thought can occur without images is older than the aphantasia literature, and an earlier episode of it deserves attention because it anticipates the present argument and also shows how the question tends to be mishandled. In the first decade of the twentieth century, experimental psychologists associated with Oswald K\"ulpe at W\"urzburg asked trained observers to perform simple acts of judgment and problem solving and then to report what had been present in consciousness. Marbe (1901), studying the judgment of weights, found that observers could not point to characteristic sensory contents on which their judgments were based, and subsequent work by Ach, Watt, Messer, and B\"uhler (the last in 1907) elaborated the finding into the claim that there are conscious contents of thinking, awarenesses of a task, a relation, or a meaning, that have no sensory or imagistic character. The claim was that thinking can be present to the thinker without being present as a picture or as inner speech. Titchener, working at Cornell, rejected the claim for years and maintained that the reports reflected a failure of introspection, the contents in question being faint kinesthetic, verbal, or visual items that trained observers would find if they looked carefully (Titchener 1909; Woodworth 1906).\footnote{For a history of the controversy see Humphrey (1951). The terms \emph{Bewusstseinslage} and \emph{Aha-Erlebnis} belong to this literature. The dispute is described here at the level needed for the argument and not in full.} The dispute was never settled by evidence. It ended when the methods that had generated it were abandoned, and its lasting effect was to demonstrate that two trained observers could examine the same task and disagree about what was in their consciousness, which weakened confidence in introspection as a court of final appeal.

The parallel with Galton is plain. In both cases a report that some thinking is without imagery met the reply that it must really contain imagery that the observer had failed to notice, and in both cases the reply was an inference from a theory and not an observation. The two episodes differ in one respect that matters. The W\"urzburg controversy was conducted among a few laboratory observers who were trained to introspect in a particular way, and it concerned whether the finding was an artifact of that training, while the later work on aphantasia concerns untrained members of the general population who describe their own experience in everyday words, which makes the sampling less artificial and the reports harder to dismiss as a product of a school.

A popular echo of the debate survives in a short text signed by L. Woolacott and dated 1920, headed \emph{Wordless Imageless Thought}, which reports a discussion among friends provoked by an article on thought by Ernest Osborne in the Australian monthly \emph{The Lone Hand}.\footnote{The text is cited from a transcription of a digital copy, reproduced as found. It is dated 1920 and was published after the June issue to which it refers; the month and the periodical of publication were not recorded. Osborne's original article was not consulted, and no biographical information about either writer was found. \emph{The Lone Hand} was a Sydney monthly of literature and poetry published from 1907 to 1928.} The piece is useful as a specimen because it shows what a general audience made of the idea within a decade or so of the laboratory debate. It begins from the possibility of thought without words or images and moves at once to a moral conclusion, that uneasy feelings before an action should be obeyed without waiting to put the reasons into words, since such knowledge of right and wrong is described as thought unclothed with words, springing from a deeper layer of the mind than that on which images and verbal expression operate. It borrows from McDougall the vocabulary of layers and levels, and it closes with the prediction that terms such as conscience, intelligence, will, and reason will pass out of use like humours and faculties, the mind being one power using the whole body.

Three features of this text are instructive. The first is the speed of the passage from a negative claim, that thought can occur without words and images, to a positive one, that there is a deeper source of knowledge from which such thought comes. The absence of an experienced picture or word is taken to indicate the presence of something more fundamental, though the absence by itself indicates nothing of the kind. In the terms of the present essay, the absence of $E$ is converted into a special kind of $R$ or of authority, without argument. The second is that the writer treats imageless and wordless as a pair, so that the picture theory and the theory of inner speech are rejected together, and the remaining thought is described only by what it lacks. The third is the prediction of vocabulary change. The writer expects the vocabulary of mental faculties to be replaced when the masses are educated, which is a recognition that the available words describe the mind in a way that research has outgrown, and which is an anticipation, from an unexpected quarter, of the claim that the vocabulary of vision should not be read as a report of what occurs. The lesson drawn here is neither the writer's nor the W\"urzburg school's. It is that the discovery that some thinking has no imagery is compatible with ordinary and unmysterious accounts of how that thinking goes, and requires no deeper level of the mind to explain it.

\section{The Picture Theory and the Grammar of Imagination}

The history of the word \emph{picture} in philosophy requires a distinction between two uses before Wittgenstein can be brought to bear on the question of imagery. In the \emph{Tractatus Logico-Philosophicus} (1921) the picture theory concerns the logical structure by which a proposition represents a possible state of affairs. A proposition is a picture in the sense that its elements stand in a determinate arrangement mirroring the arrangement of the things it depicts. This is a theory of the logical form of representation, and it is not a psychological theory.\footnote{Wittgenstein (1921). The notion of a picture there is that of a logical model of a fact, in which the form of the picture is shared with what it depicts.} It does not state that understanding a sentence requires that a person form a visual image, and the Tractatus nowhere makes such a claim. The two senses of \emph{picture}, the logical and the psychological, need to be kept apart, and confusion of them is a common source of misreading.

The later Wittgenstein addresses the psychological claim directly, and in a way that is independent of whether any given person has imagery. The \emph{Philosophical Investigations} (1953) reject the view that understanding is explained by a private mental object accompanying the words, and they do so by showing that no such object could do the explanatory work assigned to it. The pivotal passages concern the word \emph{cube}. Suppose that when the word is heard, a picture of a cube comes before the mind. The picture suggests one use of the word, but it is compatible with others, since the picture could be applied in the way of a projection that is quite different from the usual one. What fixes the application is not the picture but a way of using it, and the way of using it is not itself a further picture. The image, in other words, does not interpret itself. It stands in need of the same kind of rule-governed practice that was supposed to have been explained by its presence. A remark in the same discussion states the central point compactly: an image is not a picture, though a picture may correspond to it. The remark separates the inner episode from the thing that has a determinate use, and it denies that the former can serve as the foundation of the latter.\footnote{Wittgenstein (1953), see especially sections 139 to 141 on the cube and section 301 on images and pictures.}

The same analysis yields a distinction among capacities that tend to be conflated. There is the capacity to picture a red apple, the capacity to understand the expression \emph{red apple}, the capacity to recognize an apple when one is presented, and the capacity to use the expression appropriately in conversation. These can and usually do occur together, but they are not definitionally identical, and each can be present when another is not. A person may picture an apple vividly and misuse the word, and a person may use the word flawlessly without picturing anything. The point is not that imagery is never present or never helpful, but that it is neither necessary nor sufficient for the others.

This position had an earlier anticipation in Frege. In the distinction drawn in \emph{On Sense and Reference} (1892), the idea, or subjective mental presentation, is separated from the sense of an expression, which is objective and shared. Frege compares the referent to the moon, the sense to the real image formed by the object glass of a telescope, which can be observed by several people, and the idea to the retinal image of each observer, which belongs to one person alone.\footnote{Frege (1892). The analogy is of the moon, the real image in the telescope, and the retinal image of the observer.} In the introduction to the \emph{Foundations of Arithmetic} (1884) he lays down as a principle the separation of the psychological from the logical, and the subjective from the objective.\footnote{Frege (1884), the introduction, where the separation of the psychological from the logical is the first of three stated principles.} Different people may attach quite different images to the same word and nonetheless understand and use it in the same way, which is only possible if meaning is not identical to the image. The Fregean and Wittgensteinian arguments are not the same, since the one concerns objectivity of sense and the other concerns the application of rules, but both locate the conditions under which an expression has meaning in a shared practice rather than in any particular person's accompanying experience.

It is a further point, and an apt one for the present topic, that Wittgenstein himself diagnosed the pull of the picture. The phrase that a picture held us captive appears in connection with the conception of language as a system of names for things\footnote{Wittgenstein (1953), section 115. The remark concerns a picture of how language and logic must work, not a mental image.}, and it is a diagnosis of a philosophical habit, whereby a model of how meaning works is taken for a description of how it works. The picture of meaning as the presence of an inner object before the mind is the clearest case. It sounds natural because ordinary talk of ideas, images, and views seems to describe it. The following section considers why the talk sounds that way.

\section{The Visual Vocabulary of Ordinary Language}

Ordinary language describes understanding in visual terms with remarkable regularity. A person sees the point, has a clear idea, examines a problem from another angle, illuminates a difficulty, recognizes a pattern, envisions a future, and brings a matter into focus. A speaker may be shortsighted, may overlook a consideration, may take a narrow view, and may be told that a truth is evident, a word that descends from the Latin verb for seeing. Insight, intuition, speculation, theory, and perspective all begin in the vocabulary of sight, and a good many of the central terms of epistemology do likewise.\footnote{The cross-linguistic regularity is documented in Sweetser (1990). Examples here are drawn from English and its Latin and Greek sources.} None of these expressions entails that a visual image has occurred in the mind of the speaker or hearer, and nothing in competent use of them requires that it should.

It is useful to distinguish three phenomena that this vocabulary comprises. The first is a literal report of visual experience, as in a statement that one sees a red apple on the table or that one has a clear mental picture of the kitchen. The second is a metaphor that draws upon visual experience, in which the visual sense is active, as when a speaker deliberately describes an argument as a landscape across which the reader is to be led. The third is a conventional expression whose meaning no longer depends on any particular sensory experience, such as saying that one sees what a person means. The three shade into one another, and which is in operation on a given occasion is not always determinate, but the distinction is real, and the argument that the language of vision demonstrates the presence of internal pictures requires that most of it be of the first kind. It plainly is not. Most of it is of the third kind, and a good deal of it can be shown to be so by the fact that it is used and understood by those for whom the literal sense is unavailable, as the next section discusses.

The word \emph{idea} is especially important in this connection, and its history cautions against the assumption of a single stable meaning. It descends from a Greek word associated with seeing, and in Plato it names the form or look of a thing, the intelligible pattern apprehended by the mind rather than the eyes. Descartes used it for whatever is present to the mind in thought, including things that are not in any sense pictorial. Locke's usage was broader still, covering whatever is the object of the understanding when a person thinks. Hume restricted the term to the faint copies of impressions, so that ideas are, in his scheme, in the relevant sense images.\footnote{Hume (1739), Book I, Part I, section 1.} It is only with the last of these that the equation of an idea with a mental picture is part of the doctrine. To read the earlier uses in the same way is to project a later and narrower theory backward, and the history of the term is one of divergent usage that obscures how much of the picture-view was a theory and how much a way of talking.

Linguistic research on the semantics of perception points in the same direction. Lakoff and Johnson (1980) treat understanding as seeing as one of a family of systematic conceptual metaphors in English, and Sweetser (1990) documents a cross-linguistic regularity in Indo-European languages whereby words for physical vision come to be used for knowledge and intellectual apprehension. The regularity is historical and widely attested, and it does not, in itself, show that speakers of these languages, when they use the derived senses, experience the source domain. A community's language may preserve distinctions inherited from sensory experience without requiring that every speaker reproduce that experience internally each time the words are used. The inheritance is a feature of the lexicon and the conventions of use, while the phenomenology is a feature of individual minds, and a language can carry the first without implying a uniform distribution of the second.

The argument of this section is accordingly that the grammar of visuality is not evidence for the phenomenology of visualization. The ordinary vocabulary supports the belief that thought is picture-like for those who already hold the picture theory, because it can be read in that way. For those who do not, it is simply the vocabulary in which a long history of metaphor has been settled. The form of the vocabulary is explained by that history, and no additional hypothesis about uniform internal experience is required to explain it.

\section{Blindness, Metaphor, and the Independence of Meaning}

The philosophical argument meets unusually relevant experimental evidence in the study of language use by people who have never seen. If understanding the vocabulary of vision required the corresponding experience of vision, congenitally blind speakers would be expected to have difficulty with it. The evidence indicates that they do not, at least not in the way the strong version of the hypothesis predicts.

Landau and Gleitman (1985), in a detailed study of a blind child, showed that she acquired the meanings of \emph{look} and \emph{see} in a way that reflected their role in the language of her community and her own haptic exploration, and not in a way that reproduced the visual sense. More recently, Minervino, Martín, Tavernini, and Trench (2018) compared twenty congenitally blind and twenty sighted native speakers of Spanish on their interpretation of metaphorical expressions. In a first experiment, the expressions were of three types, visual, grasping, and other, ten of each, and the participants chose the most appropriate meaning for each; no group difference emerged for any type, including the visual. In a second experiment, expressions that were novel, with little or no occurrence in a large corpus, were used, and the blind participants understood the novel visual metaphors well in absolute terms and at levels comparable to the sighted.\footnote{Minervino et al. (2018). The stimuli were Spanish expressions, with visual items such as the Spanish equivalent of a clear book; the results concern comprehension of expressions and not the presence or absence of any inner experience.} The authors concluded that the results challenged the prediction of conceptual metaphor theory in its stronger forms, that blind people should have difficulty with metaphors whose source domain is visual.

A different line of work asks how sighted children apply visual vocabulary to blind agents. Elli, Bedny, and Landau (2021) found that children of four, six, and nine years, and adults, reliably identified the relevant body part for sighted agents, such as the eyes for \emph{look} and the ears for \emph{listen}. For blind agents, the six-year-olds, the nine-year-olds, and the adults either extended the visual verbs to other senses, as when a person sees with the hands or a cane, or said that the blind agent cannot look or see, while four-year-olds said that a blind person would use her eyes to see even as they acknowledged that her eyes do not work. The ability to apply visual verbs to agents with disabilities therefore changes between the ages of four and six, and the same pattern appeared for verbs of leg motion applied to users of wheelchairs.\footnote{The participants in Elli et al. (2021) were sighted children and adults judging blind agents; the study is evidence about how the vocabulary is learned and extended, not about blind speakers' own comprehension.} The authors take the results to indicate the importance of linguistic input and social interaction in acquiring the meanings of sensory words, and to be at odds with the view that first-person sensory experience is essential.

The evidence must be used with caution because blindness and aphantasia are different phenomena. Congenital blindness concerns the history of visual sensation, whereas aphantasia concerns the capacity for voluntary imagery in a person who may have normal sight, and neither is a surrogate for the other. A congenitally blind person may or may not have anything like visual imagery, and most people with aphantasia see normally. The two conditions are used here for a limited purpose. Each shows that a person can be a fully competent user of visual vocabulary without the experience that the vocabulary might be thought to demand, in the one case without the sensory history and in the other without the imagery. The cases support the independence of linguistic competence from the specific sensory experiences suggested by its vocabulary.

A congenitally blind speaker can therefore understand the sentence \emph{I see what you mean} without ever having had the experience of sight, and a person without voluntary visual imagery can understand an invitation to picture a better future without producing a picture. The argument is not that sensory experience plays no role in the acquisition of language, since it plainly does for most speakers and many words, but that competence with a word cannot be reduced to the sensory experiences its etymology suggests. There is an ironic parallel with Galton's discovery. Those who lack imagery have often learned about the difference late in life, and sometimes only on being asked, because the common vocabulary allowed them to take part in the practice for years without any sign that their own experience diverged from that of others.

\section{From Introspective Reports to Cognitive Mechanisms}

The historical cases raise a methodological problem, namely what reports of mental imagery can actually establish. Galton's respondents described their own experiences using words such as bright, clear, defined, and picture, and the answers depended on how those words were interpreted. Two people might report different experiences, or might report the same experience in different terms, or might apply different standards of vividness to the same inner state. The vividness questionnaire of Marks (1973) and its successors improve on the procedure by standardizing the items, but they remain self-report, and a rating of eight on a vividness scale means different things to different people. The difficulty is not peculiar to imagery; it arises wherever the object of study is a private experience that is communicated through public words. In this case it is acute because the public words were developed in a population whose members had no means of comparing notes.

Modern research has accordingly sought behavioural and physiological measures to supplement report. Keogh and Pearson (2018) used binocular rivalry, in which imagery of one stimulus biases subsequent perception, and found that participants who reported no visual imagery did not show the priming effect that imagery normally produces, which they interpreted as an absence of sensory visual imagery.\footnote{Keogh and Pearson (2018). Binocular rivalry priming is an indirect measure, and its relation to experienced vividness is itself a topic of study.} Other work has used pupillary responses and neural measures. Zeman and colleagues described, in 2010, a patient who lost the capacity for visual imagery after a medical procedure while remaining able to perform some visuospatial tasks, and in 2015 introduced the term aphantasia on the basis of people who reported never having had voluntary imagery.\footnote{Zeman et al. (2010, 2015). The 2010 case concerned a patient who lost imagery after a cardiac procedure and is reported with preserved performance on some visuospatial tasks.} A recent review in \emph{Neuropsychologia} by Scholz and colleagues (2026) argues that aphantasia challenges the view that imagery is central to cognition, and proposes an integration model on which reactivated sensory information must pass through several stages of integration before it becomes an imagery experience. On their account the neural activity observed in aphantasia is not concealed imagery but consists of sensory precursors that lack perceptual status, and these become perceptual representations on local integration and conscious imagery only on further integration with interoceptive signals.\footnote{Scholz et al. (2026). The description here follows the authors' abstract; the model is a proposal and not an established result.} Whether or not the specific model is correct, the proposal is notable for treating the difference as genuine and as occurring at a particular stage, instead of treating all reports of absent imagery as a disagreement about terminology.

The theoretical background to these developments is the long-running dispute about the format of mental representation. Shepard and Metzler (1971) showed that the time required to judge whether two figures are rotated versions of one another increases with the angle of rotation, a result that was taken to indicate analogue transformation of internal pictures, and Pylyshyn (1973) argued in reply that the data could be explained by structural descriptions without any pictorial format. The debate was never conclusively settled, and it bears on the terminology of this essay. It concerns $R$, the format in which information is held, and it can be pursued independently of $E$, whether anything is experienced as a picture. An answer to the first question does not supply an answer to the second.

Three notions should accordingly be kept distinct. The phenomenal vividness of an image is a matter of how the experience presents itself to its subject. The availability of imagistic representations is a matter of whether the system holds information in a form that supports the operations associated with imagery. The ability to perform tasks commonly associated with visualization is a matter of behaviour. Inferences between these are fallible in each direction. Successful spatial reasoning, drawing, memory, or recognition does not establish that the person has conscious visual imagery, because the same performance may be produced by other means. Conversely, the absence of reported imagery does not establish that every underlying visual or spatial representation is absent, because representation may be present without being experienced. The evidence reviewed above supports the weaker claim that conscious imagery is genuinely absent in some people, and it does not support the stronger claim that nothing imagistic is going on in them.

\section{Imagery, Imagination, and Intellectual Competence}

Returning to the historical cases, a pattern emerges that is not the one that a simple theory would predict. Galton's non-visualizers could remember the breakfast table without seeing it, and the reanalysis of his data shows that the corresponding claim about scientists was an overreach, so that the sound conclusion concerns the diversity of experience and not its correlation with intellectual calling. Swedenborg described elaborate visionary worlds in a prose that is itself systematic and analytic, and he was also a competent natural philosopher and engineer. Kant theorized imagination as a condition of cognition irrespective of what can be established about his own experience. Huxley produced richly descriptive literature while reporting unusually weak visualization. Grandin reports vivid and detailed imagery from which she derives engineering designs, and also a difficulty with algebra.

These cases invite a threefold distinction, which is the distinction with which the essay began. There is the capacity to represent something, the experience of representing it, and the ability to manipulate or communicate what is represented. Accomplishments in mathematics, literature, art, and spatial reasoning require the last of these, and probably the first, and the evidence does not show that they require the second. In his survey of mathematicians, Hadamard (1945) reported that many described their working thought as lacking vivid pictures and as depending on other kinds of signs, and he reproduced a reply from Einstein in which the elements of thought were said to be of a visual and in part muscular type, with words sought only at a secondary stage.\footnote{Hadamard (1945), Appendix II, where Einstein's reply is reproduced.} Galton too had suggested that other modes, chiefly motor, may stand in for visual imagery. Zeman and colleagues (2020), in a study of people at the extremes of imagery vividness, reported differences in the occupations most often represented among them, with scientific and mathematical callings more common among those with little or no imagery and creative callings among those with unusually vivid imagery. The finding is suggestive and is not a demonstration of causal dependence in either direction, but it is consistent with the view that different representational strategies can lead to comparable competence.

\begin{table}[ht]
\centering
\small
\begin{tabular}{p{2.6cm}p{3.7cm}p{3.7cm}p{3.2cm}}
\toprule
Case & Reported experience ($E$) & Competence ($C$) & What it shows \\
\midrule
Galton's non-visualizers & Denial of anything picture-like & Accurate recall of arrangement and content & $C$ without reported $E$ \\
Swedenborg & Waking visions of perceptual clarity, received rather than constructed & Systematic prose in theology and, earlier, in natural philosophy & Rich $E$ with uncertain voluntariness and origin \\
Grandin & Detailed, controllable imagery used as a working medium & Strong in design and spatial tasks, difficulty with algebra & Rich $E$ without general abstract competence \\
Huxley & Poor voluntary imagery & Highly descriptive writing & Visual language without corresponding $E$ \\
Congenitally blind speakers & No visual sensory history & Comparable comprehension of visual metaphors & Meaning without the sensory referent \\
Aphantasia & No voluntary visual imagery & Intact in many visuospatial tasks & $C$ and possibly $R$ without $E$ \\
\bottomrule
\end{tabular}
\caption{The cases arranged by the three components. Entries record what the sources report, and none of them is a measurement.}
\end{table}

The table is not a proof of independence, since each row rests on testimony or on limited experiments, but it shows that no one of the three components has been found to entail another. The larger philosophical implication is that cognitive competence may be realized through heterogeneous representational strategies, and that shared accomplishments do not entail identical internal experiences. This is a point that parallels, at the level of individual differences, the point made earlier about language. A community can share a practice of calculating, drawing, or describing while its members differ in what, if anything, they experience in the course of it.

\section{Conclusion: The Language of an Invisible Picture}

The essay began with Galton's breakfast table, because that request, which appeared to be a straightforward invitation to recall a familiar scene, exposed substantial differences in how individuals experienced recollection. More than a century later, ordinary language still invites the description of thought as something seen. Philosophical vocabulary, literary expression, and everyday conversation routinely employ images, perspectives, and illumination as figures for understanding, and the figures are used with confidence by speakers whose inner lives differ widely. The capacity to think about a picture is not identical to the experience of seeing one internally, and the ability to understand a visual metaphor does not depend upon having personally undergone its literal sensory referent. These two statements are instances of one point, which is that the competence displayed in a linguistic practice is not fixed by the particular experiences that individuals happen to have while exercising it.

The argument should resist two symmetrical mistakes. The first is to treat mental imagery as universally necessary for thought, a position that the evidence of divergent reports, of intact performance without reported imagery, and of competent use of visual language by the congenitally blind makes untenable. The second is to dismiss imagery as a mere linguistic illusion because some individuals do not experience it, a position that fails to respect the testimony of those who plainly do, from Swedenborg's visions to Grandin's engineering, and that would require an explanation of why their reports, unlike the reports of those who deny imagery, should be set aside. Both mistakes rest on the assumption that a single answer must hold for everyone, and Galton's observation, shorn of his theory, is that no such answer is available. The testimony of the vivid imager and the testimony of the non-visualizer can both be accepted as accurate descriptions of the respective experiences.

Several matters have been left open. The relation between reported absence of imagery and the underlying representations is under active investigation, and the proposals now in the literature may be revised. The historical cases are testimony and not data, and the two conjectural ones, that Swedenborg's experiences were of a kind captured by the modern term hyperphantasia and that Kant lacked imagery, cannot at present be confirmed. The relation between the vocabulary of vision and the history of the languages in which it occurs has been sketched only in outline. None of these gaps affects the central claim, which can be stated as a single proposition.

\begin{proposition}
The universality of a linguistic practice does not establish the universality of the phenomenal experience its vocabulary appears to describe.
\end{proposition}

The proposition is compatible with the existence of rich and vivid imagery in many people and with its absence in others. It asks only that the grammar of a shared language not be mistaken for a census of inner lives, and that the question of what is seen in the mind be put to people one at a time, as Galton first did, before it is answered for all.

\begin{thebibliography}{99}\small
\bibitem{berkeley} Berkeley, G. (1710). \emph{A Treatise Concerning the Principles of Human Knowledge}, Introduction.
\bibitem{blajenkova} Blajenkova, O., Kozhevnikov, M., and Motes, M. A. (2006). Object-spatial imagery: A new self-report imagery questionnaire. \emph{Applied Cognitive Psychology} 20, 239--263.
\bibitem{blake} Blake, W. (1790). \emph{The Marriage of Heaven and Hell}.
\bibitem{brewer} Brewer, W. F., and Schommer-Aikins, M. (2006). Scientists are not deficient in mental imagery: Galton revised. \emph{Review of General Psychology} 10, 130--146.
\bibitem{descartes} Descartes, R. (1641). \emph{Meditations on First Philosophy}, Sixth Meditation.
\bibitem{elli} Elli, G. V., Bedny, M., and Landau, B. (2021). How does a blind person see? Developmental change in applying visual verbs to blind agents. \emph{Cognition} 212, 104683.
\bibitem{frege84} Frege, G. (1884). \emph{Die Grundlagen der Arithmetik}.
\bibitem{frege92} Frege, G. (1892). \"Uber Sinn und Bedeutung. \emph{Zeitschrift f\"ur Philosophie und philosophische Kritik} 100, 25--50.
\bibitem{galton} Galton, F. (1880). Statistics of mental imagery. \emph{Mind} 5, 301--318.
\bibitem{grandin} Grandin, T. (1995). \emph{Thinking in Pictures: And Other Reports from My Life with Autism}. Doubleday. Expanded edition 2006, Vintage.
\bibitem{hadamard} Hadamard, J. (1945). \emph{The Psychology of Invention in the Mathematical Field}. Princeton University Press.
\bibitem{huxley54} Huxley, A. (1954). \emph{The Doors of Perception}. Chatto and Windus.
\bibitem{huxley56} Huxley, A. (1956). \emph{Heaven and Hell}. Chatto and Windus.
\bibitem{hume} Hume, D. (1739). \emph{A Treatise of Human Nature}, Book I, Part I.
\bibitem{kant81} Kant, I. (1781/1787). \emph{Kritik der reinen Vernunft}, chapter on the schematism of the pure concepts of the understanding.
\bibitem{kant66} Kant, I. (1766). \emph{Tr\"aume eines Geistersehers, erl\"autert durch Tr\"aume der Metaphysik}, Part I, chapter 3.
\bibitem{keogh} Keogh, R., and Pearson, J. (2018). The blind mind: No sensory visual imagery in aphantasia. \emph{Cortex} 105, 53--60.
\bibitem{koz} Kozhevnikov, M., Kosslyn, S., and Shephard, J. (2005). Spatial versus object visualizers: A new characterization of visual cognitive style. \emph{Memory and Cognition} 33, 710--726.
\bibitem{lakoff} Lakoff, G., and Johnson, M. (1980). \emph{Metaphors We Live By}. University of Chicago Press.
\bibitem{landau} Landau, B., and Gleitman, L. R. (1985). \emph{Language and Experience: Evidence from the Blind Child}. Harvard University Press.
\bibitem{locke} Locke, J. (1689). \emph{An Essay Concerning Human Understanding}, Book IV.
\bibitem{marks} Marks, D. F. (1973). Visual imagery differences in the recall of pictures. \emph{British Journal of Psychology} 64, 17--24.
\bibitem{minervino} Minervino, R. A., Mart\'in, A., Tavernini, L. M., and Trench, M. (2018). The understanding of visual metaphors by the congenitally blind. \emph{Frontiers in Psychology} 9, 1242.
\bibitem{pylyshyn} Pylyshyn, Z. W. (1973). What the mind's eye tells the mind's brain: A critique of mental imagery. \emph{Psychological Bulletin} 80, 1--24.
\bibitem{scholz} Scholz, C. O., Monzel, M., Kvamme, T. L., Liu, J., and Silvanto, J. (2026). An integration model of mental imagery and aphantasia: Conceptual framework, neuromechanistic pathways, and clinical implications. \emph{Neuropsychologia} 225, 109401.
\bibitem{shepard} Shepard, R. N., and Metzler, J. (1971). Mental rotation of three-dimensional objects. \emph{Science} 171, 701--703.
\bibitem{sweetser} Sweetser, E. (1990). \emph{From Etymology to Pragmatics}. Cambridge University Press.
\bibitem{swedenborg} Swedenborg, E. (1749--1756). \emph{Arcana Caelestia}; (1758). \emph{De Caelo et Ejus Mirabilibus, et de Inferno} (Heaven and Hell); journal of dreams, 1743--1744.
\bibitem{buhler} B\"uhler, K. (1907). Tatsachen und Probleme zu einer Psychologie der Denkvorg\"ange. \emph{Archiv f\"ur die gesamte Psychologie} 9, 297--365.
\bibitem{humphrey} Humphrey, G. (1951). \emph{Thinking: An Introduction to Its Experimental Psychology}. Methuen.
\bibitem{marbe} Marbe, K. (1901). \emph{Experimentell-psychologische Untersuchungen \"uber das Urteil}. Engelmann.
\bibitem{titchener} Titchener, E. B. (1909). \emph{Lectures on the Experimental Psychology of the Thought-Processes}. Macmillan.
\bibitem{woodworth} Woodworth, R. S. (1906). Imageless thought. \emph{Journal of Philosophy, Psychology and Scientific Methods} 3, 701--708.
\bibitem{woolacott} Woolacott, L. (1920). Wordless imageless thought. Published after June 1920; month and periodical not recorded. Responds to E. Osborne's article on thought in the June issue of \emph{The Lone Hand}.
\bibitem{witt21} Wittgenstein, L. (1921). \emph{Tractatus Logico-Philosophicus}.
\bibitem{witt53} Wittgenstein, L. (1953). \emph{Philosophical Investigations}.
\bibitem{zeman10} Zeman, A., Della Sala, S., Torrens, L. A., Gountouna, V.-E., McGonigle, D. J., and Logie, R. H. (2010). Loss of imagery phenomenology with intact visuo-spatial task performance: A case of `blind imagination'. \emph{Neuropsychologia} 48, 145--155.
\bibitem{zeman15} Zeman, A., Dewar, M., and Della Sala, S. (2015). Lives without imagery: Congenital aphantasia. \emph{Cortex} 73, 378--380.
\bibitem{zeman20} Zeman, A., et al. (2020). Phantasia: The psychological significance of lifelong visual imagery vividness extremes. \emph{Cortex} 130, 426--440.
\end{thebibliography}

\end{document}
